\chapter{Talsystemer}
\label{chap:ch2}

Vi er så vant til at arbejde i titalssystemet, at vi ofte glemmer, at det er en forholdsvis ny opfindelse, som engang blev betragtet som revolutionerende. Det er på tide at se nøje på, hvordan vi repræsenterer tal. Normalt bruger vi titalssystemet, hvor et tal som 3459 er en kort skrivemåde for \(3 \times 1000 + 4 \times 100 + 5 \times 10 + 9\). Hvert ciffers position er afgørende, fordi den gør det muligt at skelne mellem værdier som 30 og 3. Titalssystemet er et\textbf{ positionssystem}, dvs. det har faste positioner for enere, tiere, hundreder og så videre. Et ciffers position angiver den multiplikator (en potens af ti), som cifferet skal ganges med, og hver position har en værdi, der er ti gange så stor som værdien af positionen til højre for den.

Bemærk, at vi kan spare plads ved at skrive 1000 som \(10^3\), hvor eksponenten 3 angiver antallet af nuller. Altså er \(100000 = 10^5\). Hvis eksponenten er negativ, repræsenterer den en brøk, fx \(10^{-3} = \frac{1}{1000}\). Det mest geniale ved positionssystemet var måske indførelsen af decimalpunktummet, som gør det muligt at medtage decimalbrøker. For eksempel svarer tallet 123.456 til:
\[
1 \times 100 + 2 \times 10 + 3 \times 1 + 4 \times \frac{1}{10} + 5 \times \frac{1}{100} + 6 \times \frac{1}{1000}.
\]

Det kan illustreres således:
\[
\begin{array}{rccccccccc}
\text{Multiplikator:} & \ldots & 10^2 & 10^1 & 10^0 & . & 10^{-1} & 10^{-2} & 10^{-3} & \ldots \\
\text{Cifre:} & \ldots & 1 & 2 & 3 & . & 4 & 5 & 6 &\ldots \\
& & & & & \uparrow & & & & \\
& & & & & \text{Decimalpunktum:} & & & &
\end{array}
\]
Der er imidlertid ingen grund til, at vi nødvendigvis skal bruge potenser af 10, dvs. grundtallet 10. Babylonierne brugte for eksempel grundtallet 60, og grundtallet 12 var meget udbredt i middelalderens Europa. I dag er de mest anvendte talsystemer sammenfattet i \autoref{tab:number_systems}

\begin{table}[ht]
\centering
\renewcommand{\arraystretch}{1.4}
\begin{tabular}{|c|c|c|c|}
\hline
\textbf{Talsystem} & \textbf{Symboler} & \textbf{Grundtal} & \textbf{Yderligere oplysninger} \\ \hline
\textbf{Decimal} & 0-9 & 10 & - \\ \hline
\textbf{Binær} & 0, 1 & 2 & - \\ \hline
\textbf{Hexadecimal} & 0-9, A-F & 16 & $\mathrm{A} \equiv 10, \mathrm{B} \equiv 11, \mathrm{C} \equiv 12,$ $\mathrm{D} \equiv 13, \mathrm{E} \equiv 14, \mathrm{F} \equiv 15$ \\ \hline
\textbf{Oktal} & 0-7 & 8 & - \\ \hline
\end{tabular}
\caption{Oversigt over almindelige talsystemer}
\label{tab:number_systems}
\end{table}

Vi begynder med totalssystemet, som også er det system, vi behandler mest indgående i dette kapitel.

\section{Binære tal}
I totalssystemet udtrykker vi tal ved potenser af 2 i stedet for potenser af 10 som i titalssystemet. For nogle tal er det let. Husk, at $2^0=1$,

\begin{table}[!ht]
\centering
\renewcommand{\arraystretch}{1.4}
\begin{tabular}{|c c c c c c c c|}
\hline 
\multirow{2}{*}{Decimaltal} & & \multirow{2}{*}{Som potenser af 2} & \multicolumn{4}{c}{Potens af 2} & \multirow{2}{*}{Binært tal} \\
& & & 3 & 2 & 1 & 0 & \\ \hline
8 & $=$ & $2^3$ & 1 & 0 & 0 & 0 & 1000 \\
7 & $=$ & $2^2 + 2^1 + 2^0$ & 0 & 1 & 1 & 1 & 111 \\
6 & $=$ & $2^2 + 2^1$ & 0 & 1 & 1 & 0 & 110 \\
5 & $=$ & $2^2 + 2^0$ & 0 & 1 & 0 & 1 & 101 \\
4 & $=$ & $2^2$ & 0 & 1 & 0 & 0 & 100 \\
3 & $=$ & $2^1 + 2^0$ & 0 & 0 & 1 & 1 & 11 \\
2 & $=$ & $2^1$ & 0 & 0 & 1 & 0 & 10 \\
1 & $=$ & $2^0$ & 0 & 0 & 0 & 1 & 1 \\ \hline
\end{tabular}
\caption{Decimaltal i binær repræsentation}
\end{table}

Som i titalssystemet skriver vi tallet, så cifrets position angiver potensen: Den første plads efter det binære punktum er $2^{-1}$-pladsen, den næste er $2^{-2}$-pladsen og så videre. For at omregne et decimaltal til et binært tal kan vi bruge operatoren \texttt{mod}.

Se som eksempel på 88 i titalssystemet, dvs. $88_{10}$. Vi vil gerne skrive det som et binært tal. Vi tager tallet og dividerer successivt \texttt{mod} 2. Se nedenfor:

\begin{table}[h!]
\centering
\renewcommand{\arraystretch}{1.4}
\begin{tabular}{|c|c|c|c|c|}
\hline 
Trin nr. $n$ & $x_n$ & $x_n / 2$ & $x_n \bmod 2$ \\ \hline
0 & 88 & 44 & 0 \\ 
1 & 44 & 22 & 0 \\
2 & 22 & 11 & 0 \\
3 & 11 & 5 & 1 \\
4 & 5 & 2 & 1 \\
5 & 2 & 1 & 0 \\
6 & 1 & 0 & 1 \\ \hline
\end{tabular}
\caption{Omregning af decimaltallet 88 til binær form}
\end{table}

Når vi skriver den sidste kolonne i omvendt rækkefølge, dvs. nedefra og op, får vi 1011000, som er den binære form af 88, dvs. $88_{10} = 1011000_2$.

Binære tal med cifre efter punktummet er mindre almindelige, men sagtens mulige. Således er 101.1011 blot $2^2 + 2^0 + 2^{-1} + 2^{-3} + 2^{-4}$, hvilket efter lidt regning giver 5.6875. Vi har set, hvordan man omregner heltalsdelen af et decimaltal til et binært tal, og vi kan gøre det samme med en decimalbrøk. Se på 0.6875. Som før opstiller vi en tabel:

\begin{table}[h!]
\centering
\renewcommand{\arraystretch}{1.4}
\begin{tabular}{|c|c|c|c|}
\hline 
Trin nr. $n$ & $x_n$ & $x_n \times 2$ & $\left\lfloor x_n \times 2 \right\rfloor$ \\ \hline
0 & 0.6875 & 1.375 & 1 \\ \hline
1 & 0.375 & 0.75 & 0 \\ \hline
2 & 0.75 & 1.5 & 1 \\ \hline
3 & 0.5 & 1 & 1 \\ \hline
\end{tabular}
\caption{Omregning af decimalbrøken 0.6875 til binær form}
\end{table}

Når vi læser nedad, får vi $0.6875_{10} = 0.1011_2$.

\subsection*{Binær opløsning}

Fremgangsmåden fra det foregående afsnit kaldes \textbf{binær opløsning} og betegner repræsentationen af et tal i totalssystemet (grundtal 2). Ethvert decimaltal kan skrives som en sum af potenser af 2, hvor hver potens svarer til et binært ciffer (en bit) i tallets binære form.

Lad os igen se på decimaltallet 88. For at finde dets binære opløsning finder vi den største potens af 2, der er mindre end eller lig med 88, og fortsætter med at trække potenser af 2 fra, indtil vi når 0.

Først bemærker vi, at $2^6 = 64$ er den største potens af 2, der er mindre end 88:
\[
88 = 64 + 24
\]

Dernæst finder vi, at $2^4 = 16$ er den største potens af 2, der er mindre end 24:
\[
24 = 16 + 8
\]

Til sidst passer $2^3 = 8$ præcis med resten:
\[
8 = 8 + 0
\]

Altså har vi:
\[
88 = 2^6 + 2^4 + 2^3
\]

I binær form repræsenteres hver af disse potenser af 2 af et »1« på den tilsvarende plads, mens der står »0« på de pladser, hvor ingen potens af 2 bidrager:
\[
88_{10} = 1011000_2
\]

Sammenfattende:
\begin{itemize}
    \item $2^6 = 64$ svarer til det »1« længst til venstre i den binære opløsning.
    \item $2^4 = 16$ svarer til det næste »1«.
    \item $2^3 = 8$ svarer til det næste »1«.
    \item De øvrige cifre er »0«, fordi $2^5$, $2^2$, $2^1$ og $2^0$ ikke bidrager til værdien 88.
\end{itemize}

Den binære opløsning af 88 er altså $1011000_2$. Denne måde at repræsentere tal på er grundlæggende i datalogi og digital elektronik, hvor binær repræsentation er standarden for lagring og behandling af data.

\subsection*{Binære operationer}
Binære operationer er grundlæggende regneoperationer, der udføres på binære tal. Disse operationer er afgørende i computere og digitale systemer, fordi de udgør fundamentet for, hvordan computere behandler og manipulerer data.

Binær addition, subtraktion og multiplikation svarer til de tilsvarende operationer i titalssystemet, men følger enklere regler, fordi totalssystemet kun har få cifre. For eksempel følger binær addition disse regler:

\[
0 + 0 = 0, \quad 0 + 1 = 1, \quad 1 + 0 = 1, \quad 1 + 1 = 10
\]

Her giver \(1 + 1\) resultatet \(10_2\), hvilket betyder 0 med 1 i mente til den næste, højere bit. Binær subtraktion og multiplikation følger tilsvarende enkle regler, som er lette at implementere i digitale systemer.

XOR-operationen (eksklusivt OR) er en anden vigtig binær operation. XOR giver 1, hvis de to bit, der sammenlignes, er forskellige, og 0, hvis de er ens:

\[
0 \oplus 0 = 0, \quad 0 \oplus 1 = 1, \quad 1 \oplus 0 = 1, \quad 1 \oplus 1 = 0
\]

Ved binær addition bruges XOR-operationen til at lægge to bit sammen uden at tage hensyn til en eventuel mente fra en foregående bit. Det skyldes, at XOR reelt udfører addition modulo 2, hvilket passer præcis med, hvordan binær addition fungerer. For eksempel:

\[
\begin{array}{c|c|c|c}
\text{Bit 1} & \text{Bit 2} & \text{XOR (sum)} & \text{AND (mente)} \\
\hline
0 & 0 & 0 & 0 \\
0 & 1 & 1 & 0 \\
1 & 0 & 1 & 0 \\
1 & 1 & 0 & 1 \\
\end{array}
\]

I tilfældet \(1 + 1\) giver XOR summen 0, og en AND-operation (som registrerer menten) giver menten 1, hvilket resulterer i det binære tal 10.

\begin{align*}
0 + 0 &= 0 \\
0 + 1 &= 1 \\
1 + 1 &= 10 \quad \text{så vi har 1 i mente og skriver et nul} \\
1 + 1 + 1 &= 1 + (1 + 1) = 1 + 10 = 11.
\end{align*}

Vi kan opstille det på stort set samme måde som en addition i titalssystemet:

\[
\begin{array}{r|r|r|r|r|r|r|r}
& 1 & 1 & 0 & 1 & 0 & 1 & \\
+ & 1 & 0 & 1 & 1 & 1 & 0 & \\
\hline
1 & 1 & 0 & 0 & 0 & 1 & 1 & \quad \text{Sum} \\
\uparrow & & & & \uparrow & & & \\
& & & & & & &
\end{array}
\]

Den højre pil viser, hvor vi har 1 i mente. Den venstre pil viser, hvor vi har $1 + 1 + 1$, så vi har 1 i mente og 1 tilbage.

Som vi skal se nedenfor, får vi ofte brug for at håndtere flere menter. Det kan gøres på to måder, der minder om de metoder, vi kender fra titalssystemet. Vi forklarer dem med et eksempel.

\subsubsection*{Metode 1: Kolonnevis binær addition med flere menter}
Se på

\begin{equation*}
\begin{tikzpicture}[
    every node/.style={column sep=.5mm,row sep=1mm}]
    \matrix (m) [matrix of math nodes,
        nodes in empty cells,
        %nodes=draw
    ] 
    {
        &  &  &  &  &  &  &  & 1 & 1 & 1 & 1 & 1 &     \\
    +   &  &  &  &  &  &  &  & 1 & 1 & 1 & 0 & 1 &            \\
    +   &  &  &  &  &  &  &  & 1 & 1 & 1 & 0 & 1 &            \\
    +   &  &  &  &  &  &  &  & 1 & 1 & 1 & 1 & 1 &            \\
        &  &  &  &  &  &  &  &  &  &  &  &  &            \\                                                  
    };

    \draw[-,color=black,semithick] (m-4-2.south west) -- (m-4-13.south east);

\end{tikzpicture}
\end{equation*}

\textbf{Trin 1: Læg kolonnen længst til højre sammen }\newline
Begynd med at lægge bittene længst til højre sammen:

$$
1+1+1+1=100_2 \quad(\text { som er binært for } 4)
$$
Når vi læser resultatet fra højre mod venstre (dvs. fra \textit{mindst betydende bit} (LSB) til \textit{mest betydende bit}(MSB)), gør vi følgende:
\begin{itemize}
    \item skriv 0
    \item overfør 0 som mente til næste kolonne
    \item overfør 1 som mente til tredje kolonne
\end{itemize}
Du ender med

\begin{equation*}
\begin{tikzpicture}[
    row 1/.style={font=\textsl,font=\tiny, anchor=mid,
        inner sep=1.5pt},
    every node/.style={column sep=.5mm,row sep=1mm}]
    \matrix (m) [matrix of math nodes,
        nodes in empty cells,
        %nodes=draw
    ] 
    {
        &   &   &   &   &   &  &  &  &  & 1 & 0 &   &            \\
        &  &  &  &  &  &  &  & 1 & 1 & 1 & 1 & 1 &     \\
    +   &  &  &  &  &  &  &  & 1 & 1 & 1 & 0 & 1 &            \\
    +   &  &  &  &  &  &  &  & 1 & 1 & 1 & 0 & 1 &            \\
    +   &  &  &  &  &  &  &  & 1 & 1 & 1 & 1 & 1 &            \\
        &  &  &  &  &  &  &  &  &  &  &  & 0 &            \\                                                  
    };

    \draw[-,color=black,semithick] (m-5-2.south west) -- (m-5-13.south east);

\end{tikzpicture}
\end{equation*}

\textbf{Trin 2: Læg den anden kolonne fra højre sammen}\newline
Læg dernæst den anden kolonne sammen:



$$
0+1+0+0+1=10_2 \quad(\text{ som er binært for } 2)
$$
Når vi læser resultatet fra LSB til MSB:
\begin{itemize}
    \item skriv 0
    \item overfør 1 som mente til næste kolonne
\end{itemize}
Du ender med

\begin{equation*}
\begin{tikzpicture}[
    row 1/.style={font=\textsl,font=\tiny, anchor=mid,
        inner sep=1.5pt},
    row 2/.style={font=\textsl,font=\tiny, anchor=mid,
        inner sep=1.5pt},
    every node/.style={column sep=.5mm,row sep=1mm}]
    \matrix (m) [matrix of math nodes,
        nodes in empty cells,
        %nodes=draw
    ] 
    {
        &   &   &   &   &   &  &  &  &  & 1 &  &   &            \\
        &   &   &   &   &   &  &  &  &  & 1 &  &   &            \\
        &  &  &  &  &  &  &  & 1 & 1 & 1 & 1 & 1 &     \\
    +   &  &  &  &  &  &  &  & 1 & 1 & 1 & 0 & 1 &            \\
    +   &  &  &  &  &  &  &  & 1 & 1 & 1 & 0 & 1 &            \\
    +   &  &  &  &  &  &  &  & 1 & 1 & 1 & 1 & 1 &            \\
        &  &  &  &  &  &  &  &  &  &  & 0 & 0 &            \\                                                  
    };

    \draw[-,color=black,semithick] (m-6-2.south west) -- (m-6-13.south east);

\end{tikzpicture}
\end{equation*}


\textbf{Trin 3: Læg den tredje kolonne fra højre sammen}\newline
Læg nu den tredje kolonne sammen:

$$
1+1+1+1+1+1=110_2 \quad(\text{ som er binært for } 6)
$$
Når vi læser resultatet fra LSB til MSB:
\begin{itemize}
    \item skriv 0
    \item overfør 1 som mente til fjerde kolonne
    \item overfør 1 som mente til femte kolonne
\end{itemize}
Du ender med

Vores sum indtil videre:

\begin{equation*}
\begin{tikzpicture}[
    row 1/.style={font=\textsl,font=\tiny, anchor=mid,
        inner sep=1.5pt},
    every node/.style={column sep=.5mm,row sep=1mm}]
    \matrix (m) [matrix of math nodes,
        nodes in empty cells,
        %nodes=draw
    ] 
    {
        &   &   &   &   &   &  &  & 1 & 1 &  &  &   &            \\
        &  &  &  &  &  &  &  & 1 & 1 & 1 & 1 & 1 &     \\
    +   &  &  &  &  &  &  &  & 1 & 1 & 1 & 0 & 1 &            \\
    +   &  &  &  &  &  &  &  & 1 & 1 & 1 & 0 & 1 &            \\
    +   &  &  &  &  &  &  &  & 1 & 1 & 1 & 1 & 1 &            \\
        &  &  &  &  &  &  &  &  &  & 0 & 0 & 0 &            \\                                                  
    };

    \draw[-,color=black,semithick] (m-5-2.south west) -- (m-5-13.south east);

\end{tikzpicture}
\end{equation*}


\textbf{Trin 4: Læg den fjerde kolonne fra højre sammen}\newline
Gå videre til den fjerde kolonne:

$$
1+1+1+1+1=101_2 \quad(\text{ som er binært for } 5)
$$

Når vi læser resultatet fra LSB til MSB:
\begin{itemize}
    \item skriv 1
    \item overfør 0 som mente til femte kolonne
    \item overfør 1 som mente til sjette kolonne
\end{itemize}
Du ender med

\begin{equation*}
\begin{tikzpicture}[
    row 1/.style={font=\textsl,font=\tiny, anchor=mid,
        inner sep=1.5pt},
    row 2/.style={font=\textsl,font=\tiny, anchor=mid,
        inner sep=1.5pt},
    every node/.style={column sep=.5mm,row sep=1mm}]
    \matrix (m) [matrix of math nodes,
        nodes in empty cells,
        %nodes=draw
    ] 
    {
        &   &   &   &   &   &  &  & 0 &  &  &  &   &            \\
        &   &   &   &   &   &  & 1 & 1 &  &  &  &   &            \\
        &  &  &  &  &  &  &  & 1 & 1 & 1 & 1 & 1 &     \\
    +   &  &  &  &  &  &  &  & 1 & 1 & 1 & 0 & 1 &            \\
    +   &  &  &  &  &  &  &  & 1 & 1 & 1 & 0 & 1 &            \\
    +   &  &  &  &  &  &  &  & 1 & 1 & 1 & 1 & 1 &            \\
        &  &  &  &  &  &  &  &  & 1 & 0 & 0 & 0 &            \\                                                  
    };

    \draw[-,color=black,semithick] (m-6-2.south west) -- (m-6-13.south east);

\end{tikzpicture}
\end{equation*}

\textbf{Trin 5: Læg kolonnen længst til venstre sammen}\newline
Læg kolonnen længst til venstre sammen:

$$
1+1+1+1+1=101_2 \quad(\text{ som er binært for } 5)
$$
Når vi læser resultatet fra LSB til MSB:
\begin{itemize}
    \item skriv 1
    \item overfør 0 som mente til sjette kolonne
    \item overfør 1 som mente til syvende kolonne
\end{itemize}
Det giver

\begin{equation*}
\begin{tikzpicture}[
    row 1/.style={font=\textsl,font=\tiny, anchor=mid,
        inner sep=1.5pt},
    row 2/.style={font=\textsl,font=\tiny, anchor=mid,
        inner sep=1.5pt},
    every node/.style={column sep=.5mm,row sep=1mm}]
    \matrix (m) [matrix of math nodes,
        nodes in empty cells,
        %nodes=draw
    ] 
    {
        &   &   &   &   &   &  & 0 &  &  &  &  &   &            \\
        &   &   &   &   &   & 1 & 1 &  &  &  &  &   &            \\
        &  &  &  &  &  &  &  & 1 & 1 & 1 & 1 & 1 &     \\
    +   &  &  &  &  &  &  &  & 1 & 1 & 1 & 0 & 1 &            \\
    +   &  &  &  &  &  &  &  & 1 & 1 & 1 & 0 & 1 &            \\
    +   &  &  &  &  &  &  &  & 1 & 1 & 1 & 1 & 1 &            \\
        &  &  &  &  &  &  &  & 1 & 1 & 0 & 0 & 0 &            \\                                                  
    };

    \draw[-,color=black,semithick] (m-6-2.south west) -- (m-6-13.south east);

\end{tikzpicture}
\end{equation*}


\textbf{Trin 6: Læg de resterende menter sammen}\newline
Læg til sidst de resterende menter sammen:

\begin{equation*}
\begin{tikzpicture}[
    every node/.style={column sep=.5mm,row sep=1mm}]
    \matrix (m) [matrix of math nodes,
        nodes in empty cells,
        %nodes=draw
    ] 
    {
        &  &  &  &  &  &  &  & 1 & 1 & 1 & 1 & 1 &     \\
    +   &  &  &  &  &  &  &  & 1 & 1 & 1 & 0 & 1 &            \\
    +   &  &  &  &  &  &  &  & 1 & 1 & 1 & 0 & 1 &            \\
    +   &  &  &  &  &  &  &  & 1 & 1 & 1 & 1 & 1 &            \\
        &  &  &  &  &  & 1 & 1 & 1 & 1 & 0 & 0 &  0&            \\                                                  
    };

    \draw[-,color=black,semithick] (m-4-2.south west) -- (m-4-13.south east);

\end{tikzpicture}
\end{equation*}

Det følgende eksempel viser hele processen, hvor forskellige farver skelner mellem kolonnerne og de menter, de giver anledning til. Bemærk, at de to sidste cifre i summen er sorte, fordi de ikke stammer fra en bestemt kolonne, men udelukkende fremkommer af menterne.

\begin{equation*}
\begin{tikzpicture}[
    row 1/.style={font=\textsl,font=\tiny, anchor=mid,
        inner sep=1.5pt},
    row 2/.style={font=\textsl,font=\tiny, anchor=mid,
        inner sep=1.5pt},
    row 3/.style={font=\textsl,font=\tiny, anchor=mid,
        inner sep=1.5pt},
    every node/.style={column sep=.5mm,row sep=1mm}]
    \matrix (m) [matrix of math nodes,
        nodes in empty cells,
        %nodes=draw
    ] 
    {
        &   &   &   &   &   &  \textcolor{YellowOrange}{1} &  \textcolor{YellowOrange}{0} &  &  &  &  &   &            \\
        &   &   &   &   &   &  & \textcolor{Purple}{1} & \textcolor{Purple}{0} &  & \textcolor{red}{1}  &  &   &            \\
        &   &   &   &   &   &  &  & \textcolor{Green}{1}  & \textcolor{Green}{1}  & \textcolor{blue}{1}  & \textcolor{blue}{0}  &   &            \\
        &  &  &  &  &  &  &  & \textcolor{YellowOrange}{1} & \textcolor{Purple}{1} & \textcolor{Green}{1}  & \textcolor{red}{1}  & \textcolor{blue}{1} &     \\
    +   &  &  &  &  &  &  &  &  \textcolor{YellowOrange}{1} & \textcolor{Purple}{1} & \textcolor{Green}{1}  & \textcolor{red}{0} & \textcolor{blue}{1}  &            \\
    +   &  &  &  &  &  &  &  &  \textcolor{YellowOrange}{1} & \textcolor{Purple}{1} & \textcolor{Green}{1}  & \textcolor{red}{0}  & \textcolor{blue}{1}  &            \\
    +   &  &  &  &  &  &  &  &  \textcolor{YellowOrange}{1} & \textcolor{Purple}{1} & \textcolor{Green}{1}  & \textcolor{red}{1}  & \textcolor{blue}{1}  &            \\
        &  &  &  &  &  & 1 & 1 &  \textcolor{YellowOrange}{1} & \textcolor{Purple}{1} & \textcolor{Green}{0}  & \textcolor{red}{0}  & \textcolor{blue}{1}  &            \\                                                  
    };

    \draw[-,color=black,semithick] (m-7-2.south west) -- (m-7-13.south east);

\end{tikzpicture}
\end{equation*}

\subsubsection*{Metode 2: Direkte summation og forenkling}
Vi illustrerer den anden metode med det samme eksempel. I det foregående tilfælde overførte vi selve det binære tal til de næste kolonner. I denne metode skriver vi 0, hvis summen er lige, og 1, hvis summen er ulige. For hvert multiplum af 2 i summen overfører vi 1 som mente til de næste kolonner, og vi fortsætter sådan for hver kolonne, idet menterne indgår i beregningen af den pågældende kolonne.

\textbf{Trin 1: Læg kolonnen længst til højre sammen}\newline
Læg bittene i kolonne 1 sammen (talt fra LSB):

$$
1+1+1+1=100_2 \quad(\text{ som er binært for } 4)
$$
Når vi læser resultatet fra LSB til MSB:
\begin{itemize}
    \item skriv 0
    \item overfør 1 som mente for det første multiplum af 2
    \item overfør 1 som mente for det andet multiplum af 2
\end{itemize}
Det giver


\begin{equation*}
\begin{tikzpicture}[
    row 1/.style={font=\textsl,font=\tiny, anchor=mid,
        inner sep=1.5pt},
    row 2/.style={font=\textsl,font=\tiny, anchor=mid,
        inner sep=1.5pt},
    row 3/.style={font=\textsl,font=\tiny, anchor=mid,
        inner sep=1.5pt},
    every node/.style={column sep=.5mm,row sep=1mm}]
    \matrix (m) [matrix of math nodes,
        nodes in empty cells,
        %nodes=draw
    ] 
    {
        &   &   &   &   &   &  &  &  &  &  &  &   &            \\
        &   &   &   &   &   &  &  &  &  &  & 1 &   &            \\
        &   &   &   &   &   &  &  & &  &  & 1 &   &            \\
        &  &  &  &  &  &  &  & 1 & 1 & 1 & 1 & 1 &     \\
    +   &  &  &  &  &  &  &  & 1 & 1 & 1 & 0 & 1 &            \\
    +   &  &  &  &  &  &  &  & 1 & 1 & 1 & 0 & 1 &            \\
    +   &  &  &  &  &  &  &  & 1 & 1 & 1 & 1 & 1 &            \\
        &  &  &  &  &  &  &  &  &  &  &  & 0 &            \\                                                  
    };

    \draw[-,color=black,semithick] (m-7-2.south west) -- (m-7-13.south east);

\end{tikzpicture}
\end{equation*}

\textbf{Trin 2: Læg den anden kolonne fra højre sammen}\newline
Læg bittene i kolonne 2 sammen (talt fra LSB):

$$
1+1+1+1=100_2 \quad(\text{ som er binært for } 4)
$$
Når vi læser resultatet fra LSB til MSB:
\begin{itemize}
    \item skriv 0
    \item overfør 1 som mente for det første multiplum af 2
    \item overfør 1 som mente for det andet multiplum af 2
\end{itemize}
Det giver


\begin{equation*}
\begin{tikzpicture}[
    row 1/.style={font=\textsl,font=\tiny, anchor=mid,
        inner sep=1.5pt},
    row 2/.style={font=\textsl,font=\tiny, anchor=mid,
        inner sep=1.5pt},
    row 3/.style={font=\textsl,font=\tiny, anchor=mid,
        inner sep=1.5pt},
    every node/.style={column sep=.5mm,row sep=1mm}]
    \matrix (m) [matrix of math nodes,
        nodes in empty cells,
        %nodes=draw
    ] 
    {
        &   &   &   &   &   &  &  &  &  &  &  &   &            \\
        &   &   &   &   &   &  &  &  &  & 1  &  &   &            \\
        &   &   &   &   &   &  &  & &  & 1 &  &   &            \\
        &  &  &  &  &  &  &  & 1 & 1 & 1 & 1 & 1 &     \\
    +   &  &  &  &  &  &  &  & 1 & 1 & 1 & 0 & 1 &            \\
    +   &  &  &  &  &  &  &  & 1 & 1 & 1 & 0 & 1 &            \\
    +   &  &  &  &  &  &  &  & 1 & 1 & 1 & 1 & 1 &            \\
        &  &  &  &  &  &  &  &  &  &  & 0 & 0 &            \\                                                  
    };

    \draw[-,color=black,semithick] (m-7-2.south west) -- (m-7-13.south east);

\end{tikzpicture}
\end{equation*}

\textbf{Trin 3: Læg den tredje kolonne fra højre sammen}
Læg bittene i kolonne 3 sammen (talt fra LSB):

$$
1+1+1+1+1+1=110_2 \quad(\text{ som er binært for } 6)
$$
Når vi læser resultatet fra LSB til MSB:
\begin{itemize}
    \item skriv 0
    \item overfør 1 som mente for det første multiplum af 2
    \item overfør 1 som mente for det andet multiplum af 2
    \item overfør 1 som mente for det tredje multiplum af 2
\end{itemize}

Det giver


\begin{equation*}
\begin{tikzpicture}[
    row 1/.style={font=\textsl,font=\tiny, anchor=mid,
        inner sep=1.5pt},
    row 2/.style={font=\textsl,font=\tiny, anchor=mid,
        inner sep=1.5pt},
    row 3/.style={font=\textsl,font=\tiny, anchor=mid,
        inner sep=1.5pt},
    every node/.style={column sep=.5mm,row sep=1mm}]
    \matrix (m) [matrix of math nodes,
        nodes in empty cells,
        %nodes=draw
    ] 
    {
        &   &   &   &   &   &  &  &  & 1 &  &  &   &            \\
        &   &   &   &   &   &  &  &  & 1 &   &  &   &            \\
        &   &   &   &   &   &  &  & & 1 &  &  &   &            \\
        &  &  &  &  &  &  &  & 1 & 1 & 1 & 1 & 1 &     \\
    +   &  &  &  &  &  &  &  & 1 & 1 & 1 & 0 & 1 &            \\
    +   &  &  &  &  &  &  &  & 1 & 1 & 1 & 0 & 1 &            \\
    +   &  &  &  &  &  &  &  & 1 & 1 & 1 & 1 & 1 &            \\
        &  &  &  &  &  &  &  &  &  & 0 & 0 & 0 &            \\                                                  
    };

    \draw[-,color=black,semithick] (m-7-2.south west) -- (m-7-13.south east);

\end{tikzpicture}
\end{equation*}

\textbf{Trin 4: Læg den fjerde kolonne fra højre sammen}\newline
Læg bittene i kolonne 4 sammen (talt fra LSB):

$$
1+1+1+1+1+1+1=111_2 \quad(\text{ som er binært for } 7)
$$
Når vi læser resultatet fra LSB til MSB:
\begin{itemize}
    \item skriv 1
    \item overfør 1 som mente for det første multiplum af 2
    \item overfør 1 som mente for det andet multiplum af 2
    \item overfør 1 som mente for det tredje multiplum af 2
\end{itemize}
Det giver


\begin{equation*}
\begin{tikzpicture}[
    row 1/.style={font=\textsl,font=\tiny, anchor=mid,
        inner sep=1.5pt},
    row 2/.style={font=\textsl,font=\tiny, anchor=mid,
        inner sep=1.5pt},
    row 3/.style={font=\textsl,font=\tiny, anchor=mid,
        inner sep=1.5pt},
    every node/.style={column sep=.5mm,row sep=1mm}]
    \matrix (m) [matrix of math nodes,
        nodes in empty cells,
        %nodes=draw
    ] 
    {
        &   &   &   &   &   &  &  & 1 &  &  &  &   &            \\
        &   &   &   &   &   &  &  & 1 &  &   &  &   &            \\
        &   &   &   &   &   &  &  & 1 &  &  &  &   &            \\
        &  &  &  &  &  &  &  & 1 & 1 & 1 & 1 & 1 &     \\
    +   &  &  &  &  &  &  &  & 1 & 1 & 1 & 0 & 1 &            \\
    +   &  &  &  &  &  &  &  & 1 & 1 & 1 & 0 & 1 &            \\
    +   &  &  &  &  &  &  &  & 1 & 1 & 1 & 1 & 1 &            \\
        &  &  &  &  &  &  &  &  & 1 & 0 & 0 & 0 &            \\                                                  
    };

    \draw[-,color=black,semithick] (m-7-2.south west) -- (m-7-13.south east);

\end{tikzpicture}
\end{equation*}


\textbf{Trin 5: Læg kolonnen længst til venstre sammen}\newline
Læg bittene i kolonnen længst til venstre sammen:

$$
1+1+1+1+1+1+1=111_2 \quad(\text{ som er binært for } 7)
$$
Når vi læser resultatet fra LSB til MSB:
\begin{itemize}
    \item skriv resultatet i binær form, dvs. 1 1 1
\end{itemize}
Det giver


\begin{equation*}
\begin{tikzpicture}[
    row 1/.style={font=\textsl,font=\tiny, anchor=mid,
        inner sep=1.5pt},
    row 2/.style={font=\textsl,font=\tiny, anchor=mid,
        inner sep=1.5pt},
    row 3/.style={font=\textsl,font=\tiny, anchor=mid,
        inner sep=1.5pt},
    every node/.style={column sep=.5mm,row sep=1mm}]
    \matrix (m) [matrix of math nodes,
        nodes in empty cells,
        %nodes=draw
    ] 
    {
        &   &   &   &   &   &  &  &  &  &  &  &   &            \\
        &   &   &   &   &   &  &  &  &  &   &  &   &            \\
        &   &   &   &   &   &  &  & &  &  &  &   &            \\
        &  &  &  &  &  &  &  & 1 & 1 & 1 & 1 & 1 &     \\
    +   &  &  &  &  &  &  &  & 1 & 1 & 1 & 0 & 1 &            \\
    +   &  &  &  &  &  &  &  & 1 & 1 & 1 & 0 & 1 &            \\
    +   &  &  &  &  &  &  &  & 1 & 1 & 1 & 1 & 1 &            \\
        &  &  &  &  &  & 1 & 1 & 1 & 1 & 0 & 0 & 0 &            \\                                                  
    };

    \draw[-,color=black,semithick] (m-7-2.south west) -- (m-7-13.south east);
\end{tikzpicture}
\end{equation*}

hvilket svarer til det resultat, vi fik ovenfor. Selv om det ikke er vist eksplicit her, fungerer subtraktion på tilsvarende måde.

Med en af disse metoder til at håndtere flere menter kan vi også gange to binære tal med hinanden.



\subsection*{Multiplikation i totalssystemet}
Multiplikation i totalssystemet er teknisk set lettere end multiplikation i titalssystemet. I binære operationer arbejder vi udelukkende med to cifre: 0 og 1. Det betyder, at både multiplikatoren\footnote{»Multiplikanden« er det tal, der skal ganges, og »multiplikatoren« er det tal, der ganges med.} og multiplikanden kun består af 0'er og 1'er. Fremgangsmåden til at finde det binære produkt svarer til almindelig multiplikation i titalssystemet. De fire mulige produkter af enkelte binære cifre er:

\[
\begin{aligned}
& 0 \times 0=0 \\
& 0 \times 1=0 \\
& 1 \times 0=0 \\
& 1 \times 1=1 \\
& 1 \times 10_2 = 10_2 \quad \left( \text{multiplikation med grundtallet } 10_2 \text{ tilføjer et 0 til sidst} \right)
\end{aligned}
\]

Det sidste trin betyder, at $101_2 \times 10_2 = 1010_2$, hvilket svarer til tilfældet i titalssystemet: $143_{10} \times 10_{10} = 1430_{10}$.

Vi illustrerer fremgangsmåden med et par eksempler.

\begin{example}


\begin{equation*}
\begin{tikzpicture}[
    every node/.style={column sep=.5mm,row sep=1mm}]
    \matrix (m) [matrix of math nodes,
        nodes in empty cells,
        %nodes=draw
    ] 
    {
        &   &  (1) &   &   &   &  &  &  &  & 1 & 0 & 0  &            \\
     \times   &   & (2)  &   &   &   &  &  &  &  &   & 1 &  1 &            \\
        &   & (3)  &   &   &   &  &  & &  & 1 & 0 & 0  &            \\
       + &  & (4) &  &  &  &  &  &  & 1 & 0 & 0 & 0 &     \\
        &  & (5) &  &  &  &  &  &  & 1 & 1 & 0 & 0 &            \\                                                  
    };

    \draw[dashed, color=black, semithick] (m-2-2.south west) -- (m-2-13.south east);
    \draw[-,color=black,semithick] (m-4-2.south west) -- (m-4-13.south east);
    \end{tikzpicture}
\end{equation*}

Trinene er:

\begin{itemize}
    \item Gang multiplikanden (linje 1) med multiplikatorens LSB (linje 2), som her er 1.
    \item Skriv resultatet i linje 3.
    \item Tilføj et 0 i linje 4 for at tage højde for forskydningen til den næste potens af 2 i multiplikatoren.
    \item Gang multiplikanden (linje 1) med multiplikatorens næste bit (linje 2), som her også er 1.
    \item Skriv dette resultat i linje 4 foran det 0, du tilføjede før.
    \item Læg til sidst værdierne i linje 3 og 4 sammen som beskrevet i det foregående afsnit for at få det endelige resultat i linje 5.
\end{itemize}
\end{example}

Vi giver yderligere to eksempler:

\begin{example}

\begin{equation*}
\begin{tikzpicture}[
    every node/.style={column sep=.5mm,row sep=1mm}]
    \matrix (m) [matrix of math nodes,
        nodes in empty cells,
        %nodes=draw
    ] 
    {
        &   &   &   &   &   &  &  &  &  & 1 & 0 & 1  &            \\
     \times   &   &   &   &   &   &  &  &  & 1 &  0 & 1 &  1 &            \\
        &   &   &   &   &   &  &  & &  & 1 & 0 & 1  &            \\
       + &  &  &  &  &  &  &  &  & 1 & 0 & 1 & 0 &     \\
       + &  &  &  &  &  &  &  & 0 & 0 & 0 & 0 & 0 &            \\               
       + &  &  &  &  &  &  & 1 & 0 & 1 & 0 & 0 & 0 &     \\
        &  &  &  &  &  &  & 1 & 1 & 0 & 1 & 1 & 1 &            \\                                        
    };

    \draw[dashed, color=black, semithick] (m-2-2.south west) -- (m-2-13.south east);
    \draw[-,color=black,semithick] (m-6-2.south west) -- (m-6-13.south east);
    \end{tikzpicture}
\end{equation*}
    
\end{example}

\begin{example}
\label{ex1.3}
\begin{equation*}
\begin{tikzpicture}[
    every node/.style={column sep=.5mm,row sep=1mm}]
    \matrix (m) [matrix of math nodes,
        nodes in empty cells,
        %nodes=draw
    ] 
    {
        &   &   &   &   &   & \textcolor{Green}{1} & \textcolor{Green}{0} & \textcolor{Green}{0} & \textcolor{Green}{1} & \textcolor{Green}{1} & \textcolor{Green}{1} & \textcolor{Green}{0}  &            \\
     \times   &   &   &   &   &   &  &  &  &  &  1 & 0 &  1 &            \\
        &   &   &   &   &   & \textcolor{Green}{1} & \textcolor{Green}{0} & \textcolor{Green}{0} & \textcolor{Green}{1} & \textcolor{Green}{1} & \textcolor{Green}{1} & \textcolor{Green}{0}   &            \\
       + &  &  &  & \textcolor{Green}{1} & \textcolor{Green}{0} & \textcolor{Green}{0} & \textcolor{Green}{1} & \textcolor{Green}{1} & \textcolor{Green}{1} & \textcolor{Green}{0} & \textcolor{Red}{0} & \textcolor{Red}{0} &            \\             
        &  &  &  & 1 & 1 & 0 & 0 & 0 & 0 & 1 & 1 & 0 &     \\                                     
    };

    \draw[dashed, color=black, semithick] (m-2-2.south west) -- (m-2-13.south east);
    \draw[-,color=black,semithick] (m-4-2.south west) -- (m-4-13.south east);
    \end{tikzpicture}
\end{equation*}
    
\end{example}

Bemærk i \autoref{ex1.3}, at vi har udeladt den række af nuller, som multiplikatorens anden værdi ville have givet, og bemærk desuden, at vi har tilføjet to 0'er, før vi gentog multiplikanden i summen.

Binær multiplikation er ligesom binær addition en central operation i computeraritmetik. Når processen deles op i overskuelige trin -- først gange de enkelte bit og derefter lægge resultaterne sammen -- bliver det tydeligt, hvor meget den ligner de multiplikationsmetoder, vi bruger i titalssystemet. Den største forskel er, at det er enkelt og effektivt at arbejde i totalssystemet, hvor kun cifrene 0 og 1 indgår.

Vi bemærker, hvordan binær multiplikation bygger på binær addition. Hvert trin, som består af forskydninger og summer, er i bund og grund gentagne additioner, justeret med potenser af to. En solid forståelse af binær addition fører derfor naturligt til en bedre forståelse af binær multiplikation og dens anvendelser.

\section{Oktal og hexadecimal}
Ottetalssystemet (oktal) er et talsystem med grundtal 8, der bruger cifrene 0 til og med 7. Det er nært beslægtet med totalssystemet, som har grundtal 2. Forbindelsen mellem de to består i, hvor let binære tal kan omregnes til oktale tal og omvendt. Hvert oktalt ciffer svarer til præcis tre binære cifre (bit), hvilket gør omregninger ligetil. For eksempel kan det binære tal »110« omregnes direkte til det oktale ciffer »6«. På grund af denne nære sammenhæng bruges oktal ofte som en kortform for binær i datalogi, især i sammenhænge, hvor det bliver lettere at læse og fortolke binære data, når de binære cifre grupperes tre og tre.

\[
\begin{aligned}
12_8 &= 1 \cdot 8^1 + 2 \cdot 8^0  = 10_{10} \\
3021_8 & = 3 \cdot 8^3 + 0 \cdot 8^2 + 2 \cdot 8^1 + 1 \cdot 8^0 = 1553_{10}
\end{aligned}
\]

Da \(8\) er \(2^3\), kan vi udtrykke det i binær form:

\[
\begin{aligned}
3 & \rightarrow \textcolor{Green}{011} \\
0 & \rightarrow \textcolor{YellowOrange}{000} \\
2 & \rightarrow \textcolor{blue}{010} \\
1 & \rightarrow \textcolor{red}{001} \\
\end{aligned}
\]

\[
\text{Altså er } 3021_8 = \textcolor{Green}{011}\textcolor{YellowOrange}{000} \textcolor{blue}{010}\textcolor{red}{001}_2 = \textcolor{Green}{11}\textcolor{YellowOrange}{000} \textcolor{blue}{010}\textcolor{red}{001}_2
\]

Vi får det endelige resultat ved at fjerne de foranstillede nuller.

Sekstentalssystemet (hexadecimal) er et talsystem med grundtal 16, der bruger seksten forskellige symboler: cifrene 0-9 og bogstaverne A-F, hvor A repræsenterer 10, B repræsenterer 11 og så videre op til F, som repræsenterer 15. Hexadecimal er nært beslægtet med binær, fordi hvert hexadecimalt ciffer svarer præcis til fire binære cifre (bit). Denne direkte sammenhæng gør det let at omregne mellem de to systemer. For eksempel svarer det hexadecimale ciffer »A« til den binære sekvens »1010«. Fordi grupperingen er så effektiv, bruges hexadecimal ofte i datalogi som en mere kompakt og læsevenlig måde at repræsentere binære data på, især til fx hukommelsesadresser og farvekoder i webdesign.

\[
\begin{aligned}
123_{16} &= 1 \cdot 16^2 + 2 \cdot 16^1 + 3 \cdot 16^0  = 256 + 32 + 3 = 291_{10} \\
A2E_{16} & = 10 \cdot 16^2 + 2 \cdot 16^1 + 14 \cdot 16^0 = 2560 + 32 + 14 = 2606_{10}
\end{aligned}
\]

Da \(16\) er \(2^4\), kan vi for eksempel udtrykke $A2E_{16}$ i binær form:

\[
\begin{aligned}
A & \rightarrow 1010 \\
2 & \rightarrow 0010 \\
E & \rightarrow 1110 \\
\end{aligned}
\]

\[
\text{Altså er } A2E_{16} = 101000101110_2
\]

Tilsvarende får vi $5EB52_{16}$ som

\[
\begin{aligned}
5 & \rightarrow 0101 \\
E & \rightarrow 1110 \\
B & \rightarrow 1011 \\
5 & \rightarrow 0101 \\
2 & \rightarrow 0010 \\
\end{aligned}
\]

\[
\text{Altså er } 5EB52_{16} = 0101\,1110\,1011\,0101\,0010_2 = 1011110101101010010_2
\]

Grupperne med fire bit gør omregningen gennemskuelig. I den endelige repræsentation kan det ene foranstillede nul udelades uden at ændre værdien.

\section{Omregning mellem talsystemer}

Det er vigtigt i datalogi og digital elektronik at forstå, hvordan man omregner tal mellem totalssystemet, titalssystemet, ottetalssystemet og sekstentalssystemet. Hvert system har sit eget grundtal og sine egne anvendelser. Her er en trinvis vejledning, der hjælper dig med at omregne tal fra ét system til et andet.

\subsection*{Fra decimal til binær}
Sådan omregner du et decimaltal til binær form:
\begin{enumerate}
    \item Dividér decimaltallet med 2.
    \item Notér resten (den er 0 eller 1).
    \item Dividér kvotienten med 2, og notér resten.
    \item Gentag, indtil kvotienten er 0.
    \item Det binære tal er rækken af rester læst nedefra og op.
\end{enumerate}

\begin{example} Omregn \(23_{10}\) til binær form.

\begin{solution}

\[
\begin{aligned}
23 \div 2 & = 11 \quad \text{rest } 1 \\
11 \div 2 & = 5 \quad \text{rest } 1 \\
5 \div 2 & = 2 \quad \text{rest } 1 \\
2 \div 2 & = 1 \quad \text{rest } 0 \\
1 \div 2 & = 0 \quad \text{rest } 1 \\
\end{aligned}
\]
Altså er \(23_{10} = 10111_2\). \end{solution}
\end{example}

\subsection*{Fra decimal til oktal}
Sådan omregner du et decimaltal til oktal form:
\begin{enumerate}
    \item Dividér decimaltallet med 8.
    \item Notér resten.
    \item Dividér kvotienten med 8, og notér resten.
    \item Gentag, indtil kvotienten er 0.
    \item Det oktale tal er rækken af rester læst nedefra og op.
\end{enumerate}

\begin{example}Omregn \(78_{10}\) til oktal form.

\begin{solution}
    

\[
\begin{aligned}
78 \div 8 & = 9 \quad \text{rest } 6 \\
9 \div 8 & = 1 \quad \text{rest } 1 \\
1 \div 8 & = 0 \quad \text{rest } 1 \\
\end{aligned}
\]
Altså er \(78_{10} = 116_8\). \end{solution}
    
\end{example}

\subsection*{Fra decimal til hexadecimal}
Sådan omregner du et decimaltal til hexadecimal form:
\begin{enumerate}
    \item Dividér decimaltallet med 16.
    \item Notér resten (brug A, B, C, D, E, F for henholdsvis resterne 10, 11, 12, 13, 14, 15).
    \item Dividér kvotienten med 16, og notér resten.
    \item Gentag, indtil kvotienten er 0.
    \item Det hexadecimale tal er rækken af rester læst nedefra og op.
\end{enumerate}

\begin{example}Omregn \(255_{10}\) til hexadecimal form.

\begin{solution}
    
\[
\begin{aligned}
255 \div 16 & = 15 \quad \text{rest } 15 \quad (\text{F}) \\
15 \div 16 & = 0 \quad \text{rest } 15 \quad (\text{F}) \\
\end{aligned}
\]
Altså er \(255_{10} = FF_{16}\).\end{solution}

\end{example}


\subsection*{Fra binær til decimal}
Sådan omregner du et binært tal til decimal form:
\begin{enumerate}
    \item Gang hver bit med 2 opløftet til bittens position, idet du starter med 0 længst til højre.
    \item Læg alle produkterne sammen.
\end{enumerate}

Bemærk, at dette svarer til metoden med binær opløsning, som er beskrevet ovenfor.

\begin{example} Omregn \(1101_2\) til decimal form.

\begin{solution}
    

\[
\begin{aligned}
1 \cdot 2^3 + 1 \cdot 2^2 + 0 \cdot 2^1 + 1 \cdot 2^0 & = 8 + 4 + 0 + 1 = 13_{10}
\end{aligned}
\] \end{solution}

\end{example}

\subsection*{Fra binær til oktal}
Sådan omregner du et binært tal til oktal form:
\begin{enumerate}
    \item Gruppér de binære cifre tre og tre, idet du starter fra højre. Tilføj foranstillede nuller, hvis det er nødvendigt.
    \item Omregn hver gruppe med tre binære cifre til det tilsvarende oktale ciffer.
\end{enumerate}

\begin{example}Omregn \(110110_2\) til oktal form.

\begin{solution}
\[
\begin{aligned}
110 & \rightarrow 6 \\
110 & \rightarrow 6 \\
\end{aligned}
\]
Altså er \(110110_2 = 66_8\). \end{solution}

\end{example}

\subsection*{Fra binær til hexadecimal}
Sådan omregner du et binært tal til hexadecimal form:
\begin{enumerate}
    \item Gruppér de binære cifre fire og fire, idet du starter fra højre. Tilføj foranstillede nuller, hvis det er nødvendigt.
    \item Omregn hver gruppe med fire binære cifre til det tilsvarende hexadecimale ciffer.
\end{enumerate}

\begin{example}Omregn \(10110101_2\) til hexadecimal form.

\begin{solution}
\[
\begin{aligned}
1011 & \rightarrow B \\
0101 & \rightarrow 5 \\
\end{aligned}
\]
Altså er \(10110101_2 = B5_{16}\). \end{solution}

\end{example}

\subsection*{Fra oktal til binær}
Sådan omregner du et oktalt tal til binær form: Omregn hvert oktalt ciffer til det tilsvarende binære tal med 3 bit!

\begin{example}Omregn \(57_8\) til binær form.

\begin{solution}
\[
\begin{aligned}
5 & \rightarrow 101 \\
7 & \rightarrow 111 \\
\end{aligned}
\]
Altså er \(57_8 = 101111_2\). \end{solution}

\end{example}

\subsection*{Fra oktal til decimal}
Sådan omregner du et oktalt tal til decimal form:
\begin{enumerate}
    \item Gang hvert ciffer med 8 opløftet til cifrets position, idet du starter med 0 længst til højre.
    \item Læg alle produkterne sammen.
\end{enumerate}

\begin{example}Omregn \(157_8\) til decimal form.

\begin{solution}
    
\[
\begin{aligned}
1 \cdot 8^2 + 5 \cdot 8^1 + 7 \cdot 8^0 & = 64 + 40 + 7 = 111_{10}
\end{aligned}
\] \end{solution}

\end{example}

\subsection*{Fra oktal til hexadecimal}
Sådan omregner du et oktalt tal til hexadecimal form:
\begin{enumerate}
    \item Omregn først det oktale tal til binær form.
    \item Omregn derefter det binære tal til hexadecimal form ved at gruppere de binære cifre fire og fire.
\end{enumerate}

\begin{example}Omregn \(157_8\) til hexadecimal form.

\begin{solution}

\[
\begin{aligned}
1 & \rightarrow 001 \\
5 & \rightarrow 101 \\
7 & \rightarrow 111 \\
\end{aligned}
\]
Altså er \(157_8 = 001101111_2 = 6F_{16}\). \end{solution}
\end{example}

\subsection*{Fra hexadecimal til binær}
Sådan omregner du et hexadecimalt tal til binær form: Omregn hvert hexadecimalt ciffer til det tilsvarende binære tal med 4 bit.

\begin{example}Omregn \(2B_{16}\) til binær form.

\begin{solution}
    


\[
\begin{aligned}
2 & \rightarrow 0010 \\
B & \rightarrow 1011 \\
\end{aligned}
\]
Altså er \(2B_{16} = 00101011_2\).\end{solution}

\end{example}

\subsection*{Fra hexadecimal til decimal}
Sådan omregner du et hexadecimalt tal til decimal form:
\begin{enumerate}
    \item Gang hvert ciffer med 16 opløftet til cifrets position, idet du starter med 0 længst til højre.
    \item Læg alle produkterne sammen.
\end{enumerate}

\begin{example}Omregn \(2B_{16}\) til decimal form.

\begin{solution}
\[
\begin{aligned}
2 \cdot 16^1 + 11 \cdot 16^0 & = 32 + 11 = 43_{10}
\end{aligned}
\] \end{solution}
\end{example}

\subsection*{Fra hexadecimal til oktal}
Sådan omregner du et hexadecimalt tal til oktal form:
\begin{enumerate}
    \item Omregn først det hexadecimale tal til binær form.
    \item Omregn derefter det binære tal til oktal form ved at gruppere de binære cifre tre og tre.
\end{enumerate}

\begin{example}Omregn \(2B_{16}\) til oktal form.

\begin{solution}

\[
\begin{aligned}
2 & \rightarrow 0010 \\
B & \rightarrow 1011 \\
\end{aligned}
\]
Altså er \(2B_{16} = 00101011_2 = 53_8\). \end{solution}

\end{example}

\subsection*{Afsluttende bemærkninger om omregning}
Begrebet opløsning spiller en central rolle i disse omregninger. Uanset om du opløser et decimaltal i dets binære, oktale eller hexadecimale form eller omregner et binært tal til dets oktale eller hexadecimale modstykke, udtrykker du mere eller mindre tallet ved potenser af grundtallet. Metoden med opløsning er i bund og grund den samme for alle systemer, da den går ud på rekursivt at dividere med den højeste potens af grundtallet:

$$
\begin{aligned}
7562_{10} & =1 \cdot 16^3+3466=1 \cdot 16^3+13 \cdot 16^2+138 \\
& =\textcolor{red}{1} \cdot 16^3+ \textcolor{blue}{13} \cdot 16^2+\textcolor{YellowOrange}{8} \cdot 16^1+\textcolor{Green}{10} \cdot 16^0=\textcolor{red}{1} \, \textcolor{blue}{D} \, \textcolor{YellowOrange}{8} \, \textcolor{Green}{A}
\end{aligned}
$$

Når du forstår disse opløsninger og sammenhængene mellem talsystemerne, kan du effektivt skifte mellem dem, så du kan repræsentere og behandle data i det format, der passer bedst til den givne situation.
